\documentclass[11pt,a4paper]{article}
\usepackage{gtu-el-yazisi}
\begin{document}
\notbilgi{Sayısal İşaret İşleme}{Ayrık zamanlı sistemler, LTI, konvolüsyon, FIR/IIR}{Dsp-ders-2.pdf}

\sayfa{1}
\[ x(t)=\cos(2\pi t) \]
$F_s = 5\,\mathrm{Hz}$ olsun
\[ x[n] = x(t)\Big|_{t=nT_s} = \cos\frac{2\pi n}{5} = \cos 0{,}4\pi n \]

\[ \omega = 0{,}4\pi \]
\[ 2\pi k = 0{,}4\pi N \]
\[ N = 5k \qquad N_0 = 5 \ \text{periyottur.} \]

\[ x_1[n] = \cos(0{,}4\pi n) \qquad \leftarrow \text{bu işaretin frekansı daha yüksek} \]
\[ x_2[n] = \cos(0{,}2\pi n) \]

\begin{dikkat}
Daha sık örnekleme aynı işaret için ayrık zamanlı işaretin frekansını azaltır
\end{dikkat}

Ayrık zamanda daha yüksek değişimler daha yüksek frekans demektir, az değişler daha az frekans demektir.

Ayrık zamanda en yüksek frekans $\omega=\pi$'dir
\[ x[n]=\cos\pi n \]
\elcizimi[0.35]{sekiller/Dsp-ders-2-s1-1.png}[1]

\section{Systems}
\elcizimi[0.4]{sekiller/Dsp-ders-2-s1-2.png}[1]
\[ y[n] = x[n-n_d] \qquad \leftarrow \text{gecikme} \]

\textit{bellekli} $\swarrow$
\[ y[n] = \tfrac{1}{3}x[n] + \tfrac{1}{3}x[n-1] + \tfrac{1}{3}x[n-2] \]
ortalama alıyor, hareketli (moving average) (son üç örneklemin ortalaması alıyor)

\elcizimi[0.7]{sekiller/Dsp-ders-2-s1-3.png}[1]
(Alçak geçiren filtre aslında ayrık zamanda)

\sayfa{2}
\[ y[n]=\frac{1}{4}x[n] \qquad \text{belleksiz} \]

\section{Linear system}
linear system
\elcizimi[0.4]{sekiller/Dsp-ders-2-s2-1.png}[2]
\elcizimi[0.55]{sekiller/Dsp-ders-2-s2-2.png}[2]
\elcizimi[0.85]{sekiller/Dsp-ders-2-s2-3.png}[2]

\[ y[n]=8x[n]+1 \]
linear systeme 0 verince 0 çıkmalıdır
\begin{align*}
x_1[n]&=2 & y_1[n]&=17\\
x_2[n]&=3 & y_2[n]&=25\\
x_3[n]&=5 & y_3[n]&=41\neq 42
\end{align*}

\section{Time invariant Systems}
\[ x_1[n]\rightarrow y_1[n] \quad \text{time invariant ise} \]
\[ x_1[n-n_d]\rightarrow y_1[n-n_d] \]
\elcizimi[0.8]{sekiller/Dsp-ders-2-s2-4.png}[2]

\sayfa{3}
\section{Nedensel sistemler}
çıkışı şimdiki ve önceki örnekleme bağlı

zamansa nedensel olmayan sistemler gerçekleştirilemez\\
Ama konumsa gerçekleşebilir!!

non-causal -- $y[n]=x[n-1]+x[n]+x[n+1]$\\
anti causal -- $y[n]=x[n+1]$ \quad sadece geleceğe bağlı

\section{Kararlılık (BI-BO)}
\elcizimi[0.6]{sekiller/Dsp-ders-2-s3-1.png}[3]
\begin{align*}
|x[n]| &< B_x \le \infty\\
|y[n]| &< B_y \le \infty
\end{align*}

Bir sistemin kararlı olduğunu nasıl anlarız\\
kararlı olduğunu göstermek çok zor

\[ y[n]=(x[n])^3 \]
\begin{align*}
|x[n]| &< B_x\\
|y[n]| &< (B_x)^3 \quad \rightarrow \text{hala sonlu}
\end{align*}

\[ y[n]=\begin{cases} 0, & n<0\\ n+1, & n\ge 0 \end{cases} \]
$\uparrow$ $n$ sonsuza giderken sistem kararsızdır.

\sayfa{4}
\section{LTI Systems}
\elcizimi[0.5]{sekiller/Dsp-ders-2-s4-1.png}[4]
\elcizimi[0.35]{sekiller/Dsp-ders-2-s4-2.png}[4]
\[ y[n]=\sum_{k=-\infty}^{\infty} x[k]\,h[n-k] = x[n]*h[n] \]
\textit{convolution sum}

\elcizimi[0.6]{sekiller/Dsp-ders-2-s4-3.png}[4]
\begin{align*}
y[n] &= \sum_{k=-\infty}^{\infty} x[k]\,h[n-k]\\
&= \sum_{k=-\infty}^{\infty} h[k]\,x[n-k]
\end{align*}

\begin{align*}
x[n] &= \begin{bmatrix} 1 & 2 & 3 \end{bmatrix}\\
h[n] &= \begin{bmatrix} 2 & 1 & 3 \end{bmatrix}
\end{align*}

\sayfa{5}
\elcizimi[0.7]{sekiller/Dsp-ders-2-s5-1.png}[5]
\elcizimi[0.7]{sekiller/Dsp-ders-2-s5-2.png}[5]
\[ y[0]=2\cdot 1=2 \]
\elcizimi[0.5]{sekiller/Dsp-ders-2-s5-3.png}[5]
\[ y[1]=1\cdot 1+2\cdot 2=5 \]
$k=0 \qquad 1$

\sayfa{6}
\begin{tabular}{@{}l@{\qquad}l@{\qquad}l@{}}
$x$ & boyutu 3 & $= A$\\
$h$ & boyutu 3 & $= B$\\
$y$ & boyutu 5 & $= A+B-1$
\end{tabular}

\elcizimi[0.6]{sekiller/Dsp-ders-2-s6-1.png}[6]
\begin{align*}
w[n] &= x[n]*h_1[n]\\
y[n] &= w[n]*h_2[n]\\
&= (x[n]*h_1[n])*h_2[n]\\
&= x[n]*(h_1[n]*h_2[n])
\end{align*}

\elcizimi[0.85]{sekiller/Dsp-ders-2-s6-2.png}[6]
\elcizimi[0.7]{sekiller/Dsp-ders-2-s6-3.png}[6]

\begin{align*}
h[n] \qquad y[n] &= x[n]*h[n]\\
&= \sum_{k=-\infty}^{\infty} h[k]\,x[n-k]\\
&= \sum_{k=-\infty}^{\infty} x[k]\,h[n-k]
\end{align*}
\begin{align*}
y[n] &= \cdots + h[-2]\,x[n-2] + h[-1]\,x[n-1]\footnotemark\\
&\quad + h[0]\,x[n] + h[1]\,x[n-1] + h[2]\,x[n-2] + \cdots
\end{align*}
\footnotetext{Claude notu: muhtemelen $x[n+2]$ ve $x[n+1]$ olmalı.}

\sayfa{7}
Nedensel sistemlerde çıkış gelecekteki girişlerden bağımsız

\begin{align*}
h[-1] &= 0\\
h[-2] &= 0\\
&\ \ \vdots\\
h[-n] &= 0
\end{align*}
\[ h[n]=0 \quad n<0 \]
$\Rightarrow$ sistem nedenseldir.

\vurgu[pembe]{Kararlılık} = Genliği sınırlı bir giriş için çıkışı sınırlı olan sistem
\[ \sum_{n=-\infty}^{\infty} |h[n]| < \infty \qquad \text{(mutlak toplanabilir)} \]

\begin{dikkat}
Sakın sınavda mutlak koymayı unutma 0 alırsın :)
\end{dikkat}

\[ h[n]=u[n] \qquad \text{kararsız} \]
\[ h[n]=\alpha^{n}u[n] \]
geometrik seri toplamı
\formulkutu[önemli]{\text{toplam: } \frac{1}{1-|\alpha|} \qquad |\alpha|<1}
\[ |\alpha|\ge 1 \ \text{ise kararsızdır sistem} \]

\[ h[n]=0{,}2^{n}u[n] \qquad \text{toplam } \frac{1}{1-|\alpha|}=\frac{1}{1-0{,}2}=1{,}25 \]
\[ = 1+0{,}2+0{,}04+\cdots \]

\section{FIR ve IIR}
Finite Impulse Response\\
Infinite Impulse Response

\sayfa{8}
\elcizimi[0.85]{sekiller/Dsp-ders-2-s8-1.png}[8]
\elcizimi[0.9]{sekiller/Dsp-ders-2-s8-2.png}[8]
\elcizimi[0.8]{sekiller/Dsp-ders-2-s8-3.png}[8]

$h[n]$'in genliği her zaman sonludur

FIR her zaman kararlı

IIR (her zaman) kararlı değil ama kullanırız çünkü daha az pencere (örnek) ile FIR'ın yaptığı iş yapılabilir (bazı işaretler için)

\end{document}
